\appendix
\section{Summary of Results}
\label{sec:appendix}

\subsection{Task Completion Matrix}

\begin{table}[h]
\centering
\begin{tabular}{@{}llll@{}}
\toprule
\textbf{Task} & \textbf{Description} & \textbf{Result} & \textbf{Reference} \\
\midrule
1.1 & Reverse-sorted predicate & Definition~\ref{def:rsorted} & \S\ref{sec:specifications} \\
1.2 & H-index predicate & Definition~\ref{def:hindex} & \S\ref{sec:specifications} \\
\midrule
2.1 & Safety of \texttt{compute} & Theorem~\ref{thm:compute-safety} & \S\ref{sec:compute} \\
2.2 & Correctness of \texttt{compute} & Theorem~\ref{thm:compute-correct} & \S\ref{sec:compute} \\
\midrule
3.1 & Safety of \texttt{compute\_opt} & Theorem~\ref{thm:opt-safety} & \S\ref{sec:compute_opt} \\
3.2 & Correctness of \texttt{compute\_opt} & Theorem~\ref{thm:opt-correct} & \S\ref{sec:compute_opt} \\
\midrule
4.1 & Safety of \texttt{update} & Theorem~\ref{thm:update-safety} & \S\ref{sec:update} \\
4.2 & Sorted preservation & Theorem~\ref{thm:update-sorted} & \S\ref{sec:update} \\
4.3 & Contents consistency & Theorem~\ref{thm:update-contents} & \S\ref{sec:update} \\
4.4 & H-index correctness & Theorem~\ref{thm:update-hindex} & \S\ref{sec:update} \\
\midrule
B.1 & Complexity bounds & Theorems~\ref{thm:compute-complexity}--\ref{thm:update-complexity} & \S\ref{sec:bonus} \\
B.2 & Sort + compose & Theorem~\ref{thm:composition} & \S\ref{sec:bonus} \\
B.3 & Generalized \texttt{update\_k} & Theorem~\ref{thm:updatek} & \S\ref{sec:bonus} \\
\bottomrule
\end{tabular}
\caption{Summary of all 14 tasks (13 numbered + B.3) with theorem references.}
\label{tab:summary}
\end{table}

\subsection{Proof Dependencies}

The proof structure forms a directed acyclic graph:

\begin{center}
\begin{tabular}{@{}ll@{}}
\toprule
\textbf{Result} & \textbf{Depends on} \\
\midrule
Theorem~\ref{thm:unique} (uniqueness) & Definition~\ref{def:hindex}, Lemma~\ref{lem:countge-mono} \\
Lemma~\ref{lem:transition} (transition point) & Definitions~\ref{def:rsorted},~\ref{def:hindex},~\ref{def:countge} \\
Corollary~\ref{cor:transition-equiv} & Lemma~\ref{lem:transition} \\
Theorem~\ref{thm:compute-correct} & Corollary~\ref{cor:transition-equiv}, Invariant~\ref{inv:compute} \\
Theorem~\ref{thm:opt-correct} & Corollary~\ref{cor:transition-equiv}, Invariant~\ref{inv:compute-opt} \\
Lemma~\ref{lem:search-correct} & Invariant~\ref{inv:update-search}, Lemma~\ref{lem:equal-block} \\
Theorem~\ref{thm:update-sorted} & Lemma~\ref{lem:search-correct} \\
Theorem~\ref{thm:update-contents} & Lemma~\ref{lem:search-correct} \\
Lemma~\ref{lem:hindex-incr} & Definition~\ref{def:hindex}, Lemma~\ref{lem:countge-mono} \\
Theorem~\ref{thm:update-hindex} & Corollary~\ref{cor:transition-equiv}, Lemma~\ref{lem:hindex-incr},~\ref{lem:search-correct} \\
Theorem~\ref{thm:composition} & Theorems~\ref{thm:rsorted-correct},~\ref{thm:opt-correct} \\
Theorem~\ref{thm:updatek} & Lemma~\ref{lem:search-correct}, Theorem~\ref{thm:opt-correct} \\
\bottomrule
\end{tabular}
\end{center}

\subsection{Key Invariants at a Glance}

\begin{table}[h]
\centering
\small
\begin{tabular}{@{}lp{10cm}@{}}
\toprule
\textbf{Function} & \textbf{Loop Invariant} \\
\midrule
\texttt{compute} &
$I(h): 0 \leq h \leq n \;\wedge\; \forall\, j < h.\; a[j] \geq h$ \\
\texttt{compute\_opt} &
$J(\mathit{lo},\mathit{hi}): 0 \leq \mathit{lo} \leq \mathit{hi} \leq n
\;\wedge\; (\forall\, j < \mathit{lo}.\; a[j] > j)
\;\wedge\; (\forall\, j \geq \mathit{hi},\, j < n.\; a[j] \leq j)$ \\
\texttt{update} (search) &
$K(\mathit{lo},\mathit{hi}): 0 \leq \mathit{lo} \leq \mathit{hi} \leq i
\;\wedge\; (\forall\, j < \mathit{lo}.\; a[j] > x)
\;\wedge\; a[\mathit{hi}] = x$ \\
\bottomrule
\end{tabular}
\caption{Loop invariants used in the proofs.}
\label{tab:invariants}
\end{table}

\subsection{Discussion}

\subsubsection{The Transition-Point Characterization}

The most important insight in our verification is Lemma~\ref{lem:transition}, which
reduces the h-index from a counting definition (involving $\countge$, which requires
examining all $n$ elements) to a positional property that can be checked at a single
array position. This characterization is what makes both the linear-scan and
binary-search algorithms work: they search for the \emph{transition point}
where $a[j]$ crosses from $> j$ to $\leq j$.

For unsorted arrays, no such positional characterization exists --- one must examine
all elements to compute $\countge$. This explains why the challenge specifies
reverse-sorted input.

\subsubsection{Why Leftmost in Update}

The \texttt{update} function's binary search for the leftmost occurrence of $x$
is not arbitrary. Consider a reverse-sorted array where positions
$p, p{+}1, \ldots, q$ all have value $x$. To increment one of these values to
$x+1$ while maintaining the sorted order, we must choose position $p$ (the leftmost):
\begin{itemize}
  \item Position $p$: After $a[p] \gets x+1$, we need $a[p{-}1] \geq x+1$ (true,
    since $a[p{-}1] > x$ by the block structure) and $a[p{+}1] \leq x+1$
    (true, since $a[p{+}1] = x \leq x+1$).
  \item Any position $m > p$: After $a[m] \gets x+1$, we'd need $a[m{-}1] \geq x+1$.
    But $a[m{-}1] = x$ (since $m{-}1 \geq p$ and all such positions have value $x$).
    So $a[m{-}1] = x < x+1 = a[m]$, violating the sorted order.
\end{itemize}

\subsubsection{Correctness of the H-Index Update Logic}

The \texttt{update} function avoids recomputing the h-index from scratch.
Instead, it uses the observation that incrementing one citation can change
the h-index by at most~1 (Lemma~\ref{lem:hindex-incr}). The condition
$\mathit{lo} = h \wedge a[\mathit{lo}] = h + 1$ precisely captures when the
h-index increases: the element at the transition point (position $h$) was the
``blocking'' element with $a[h] = h \leq h$, and after incrementing to $h+1$,
it no longer blocks, pushing the transition point one position further.
